% Source: https://artofproblemsolving.com/community/c3h2007214_2017_state_sprint_27
% Stable identity: aops-topic-2007214-q02
An assignment is a function from the $20$ coins to the $4$ people. Call a person \emph{poor} if he receives $0$ or $1$ coins; we count assignments with no poor person by inclusion--exclusion.

Fix a set $T$ of $k$ people and count the assignments in which everyone in $T$ is poor. Suppose exactly $j$ of them receive one coin ($0\le j\le k$). Choose these $j$ people in $\binom kj$ ways and give them distinct coins in $20\cdot19\cdots(21-j)$ ways; the other $20-j$ coins go to the $4-k$ people outside $T$ in $(4-k)^{20-j}$ ways. Hence the count for $T$ is
\[
f(k)=\sum_{j=0}^{k}\binom kj\,\dfrac{20!}{(20-j)!}\,(4-k)^{20-j}.
\]
\begin{itemize}
\item $f(0)=4^{20}=1099511627776$.
\item $f(1)=3^{20}+20\cdot3^{19}=3486784401+23245229340=26732013741$.
\item $f(2)=2^{20}+2\cdot20\cdot2^{19}+20\cdot19\cdot2^{18}=1048576+20971520+99614720=121634816$.
\item $f(3)=1+3\cdot20+3\cdot20\cdot19+20\cdot19\cdot18=1+60+1140+6840=8041$.
\item $f(4)=0$, since four poor people can hold at most $4$ coins.
\end{itemize}
By inclusion--exclusion the number of assignments with nobody poor is
\[
f(0)-4f(1)+6f(2)-4f(3)+f(4)=1099511627776-106928054964+729808896-32164=993313349544 .
\]
The answer is $\boxed{993313349544}$.
